\section{Drupal Project Analysis} % {{{

\marginnote{\includegraphics[width=\marginparwidth]{drupal/logo}}

Drupal is a \ac{CMS} and a \ac{CMF}. It is written in the PHP programming
language and licensed under the \ac{GPL}. Drupal is used for a wide range
of capabilities such as blogs, company and political websites, communities and
is directed to be a widely spread cms \cite{DrupalOverview}. Drupal is also
known as \emph{Drupal Core} inside the Community and contains the base
functionality of the cms. However it can be extended through numerous modules.
In the following though only Drupal Core will be analyzed.

\begin{figure}[htbp]
  \centering
  \includegraphics[width=0.8\textwidth]{drupal/commits_by_author}
  \caption{Monthly activity of the most active Drupal Core developers. It
  clearly shows the involvement of the founder Dries Buytaert. In recent years
  however the roles of Gábor Hojtsy and Angie Byron became more important.}
\end{figure}

\subsection{History} % {{{

The first version of Drupal was released in 2001 by Dries Buytaert
\cite{DrupalHistory}. Originally the software provided a black board like
directory. Quickly however a fully-fledged \ac{CMS} emerged out of the initial
version. The name originates from the dutch word \emph{druppel} which
translated means \emph{drop}. Especially in the last years it
experienced a wide distribution and is used for at least \unit[1.7]{\%} of all
websites worldwide and holds a marketshare of \unit[6.2]{\%}
\cite{DrupalBuiltWith,DrupalW3Techs}.

% }}}

\subsection{Community} % {{{

The Drupal project has a large user and developer community around the world.
According to own statements, over 250.000 users were registered on
\url{drupal.org} in August 2011. About 1200 of those were additionally listed
as developers, whereas the number could be a bit higher \cite{DrupalBuytaert}.

Twice a year the community holds a conference under the name \emph{Drupal
Conference} or \emph{DrupalCon}. To alleviate the journey of the attendees the
conference takes place alternatively in Europe and North America.

The communication possibilities sprawl from forums on \url{drupal.org}, mailing
lists and discussion groups on \url{groups.drupal.org}. Additionally a wide
range of \ac{IRC} channels for different aspects of the Drupal project are
being used.

Next to the last-mentioned classification of users and developers, people can
-- provided that they are working on Drupal Core -- be classified into one of
the following categories \cite{DrupalCoreDevelopers}.

\paragraph{Founder and Lead Developer}

Dries Buytaert, the founder of the Drupal project, holds the role of the
project leader and therefore decides the project's direction. Furthermore he is
the main decision maker. Though in some cases he also gives away his decision
control to a trusted person.

\begin{table}[h!t]
  \centering
  \begin{tabularx}{\textwidth}{Xlr}
    \toprule
    \tableheadline{Venue} & \tableheadline{Date}  & \tableheadline{Attendees} \\
    \midrule
    Munich, Germany       & August 2012           & N/A \\
    Denver, USA           & March 2012            & N/A \\
    London, England       & August 2011           & 1751 \\
    Chicago, USA          & March 2011            & 3000 \\
    Kopenhagen, Denmark   & August 2010           & 1200 \\
    San Francisco, USA    & April 2010            & 3000 \\
    Paris, France         & September 2009        & 850 \\
    Washington D.C., USA  & March 2009            & 1400 \\
    Szeged, Hungary       & August 2008           & 500 \\
    Boston, USA           & March 2008            & 850 \\
    Barcelona, Spain      & September 2007        & 450 \\
    Sunnyvale, USA        & March 2007            & \textasciitilde 300 \\
    Brussels, Belgium     & September 2006        & 150 \\
    Vancouver, Canada     & February 2006         & \textasciitilde 150 \\
    Amsterdam, Netherlands& October 2005          & \textasciitilde 100 \\
    Portland, USA         & August 2005           & 100 \\
    Antwerpen, Belgium    & February 2005         & \textasciitilde 50 \\
    \bottomrule
  \end{tabularx}
  \caption[Previous and planned Drupal conferences]{Previous and planned DrupalCon events \cite{DrupalWalling}.}
\end{table}

\paragraph{Core Committer}

Only few people have write access to the Drupal Core repository. They are known
as Core Committer and look through code changes and propositions and decide
about patch acceptances. At this juncture there is a distinction between
permanent Core Committers and Branch Maintainers.

\begin{description}
  \item[Permanent Core Committer] Dries Buytaert is currently the only permanent core
    committer.
  \item[Branch Maintainer] Next to Dries Buytaert they are in charge of
  maintaining a particular Drupal version.
  \begin{itemize}
    \item Gerhard Killesreiter for Drupal~4.7.x.
    \item Neil Drumm for Drupal~5.x.
    \item Gábor Hojtsy for Drupal~6.x.
    \item Angie Byron for Drupal~7.x.
  \end{itemize}
\end{description}

\paragraph{Maintainer}

While maintainers often aren't involved in the decision making progress, they
have the liability of small portions of Drupal Core. Most of them are
particular Core modules or technical domains such as JavaScript. The assignment
happens over Dries Buytaert, whom interested developers can approach or will be
invited to join.

\paragraph{Core Contributor}

All developers who bring in code or documentation for Drupal Core are known as
Core Contributors. All modification proposals are checked by another Core
Committer in a review process and then either accepted or dropped.

\begin{figure}[htbp]
  \centering
  \includegraphics[width=0.8\textwidth]{drupal/commits_by_year}
  \caption{Yearly overview of commits to Drupal Core. The peak in 2009 and 2010
  is most likely related to the development of Drupal~7 which was released in
  early 2011.}
\end{figure}

% }}}

\subsection{Release Process} % {{{

For releases, Drupal uses a version naming with two numbers since Drupal~5. The
numbers stand for major and minor versions \cite{DrupalUpgrade}. New major
releases are quite scarce and will get released every few years. They are
planned for a long time in advance, often with incompatible changes. Before
Drupal~5 a three numbers versioning scheme was used.

While the time periods between single Drupal Core major releases are not bound,
the time until a stable release is divided into single phases
\cite{DrupalReleaseCycle}. The length of a single phase is not fixated either
but limited by a range of factors which will be explained in the following.

\begin{figure}[htbp]
  \centering
  \includegraphics[width=\textwidth]{drupal/releases}
  \caption{Major releases of Drupal.}
\end{figure}

\paragraph{Code Thaw}

After each major release a new branch will be created in the Drupal Core
repository. In this phase all new features and modifications can be discussed
and added. No restrictions apply.

\begin{figure}[htbp]
  \centering
  \includegraphics[width=\textwidth]{drupal/punchcard}
  \caption{Time based view on commits of core contributors. There seems to be a
  tendency to evenings during the week and Sundays which could be interpreted
  with a more community driven project.}
\end{figure}

\paragraph{Code Slush}

This phase was added to the Drupal~7 release cycle to restrict the time between
major releases. Important features will be elaborated and handled with the help
of so called initiatives. In this phase no new features except the chosen
initiatives are allowed. Though changes will be accepted to fix errors or add
missing functionality.

At the end of this phase, the \ac{API} freeze will be declared. At this point
in time, the \ac{API} can not be changed anymore except critical errors are
found and have to be corrected.

\begin{figure}[htbp]
  \centering
  \includegraphics[width=0.8\textwidth]{drupal/release_phases}
  \caption{Drupal release process phases.}
\end{figure}

\paragraph{Polish Phase}

In this phase Drupal Core will be polished and the last enhancements on
performance, accessibility, documentation and other are brought in. At the end
of the phase string freeze and user interface freeze will be put in place.

\paragraph{Release Phase}

Lastly there will be several alpha, beta, release candidates and in the long
run the finally released. The first release candidate will be published once
there are no more critical bugs present. Additional release candidates will be
released if new critical bugs appear, otherwise, if no new critical bugs
appear, the final release is published.

\begin{figure}[htbp]
  \centering
  \includegraphics[width=0.8\textwidth]{drupal/commits_by_month}
  \caption{Amount of commits per month of core contributors. Again, the peak in
  2009 and 2010 seems to match the development of Drupal~7. The low thereafter
  is then related to the planning phase of Drupal~8.}
\end{figure}

% }}}

\subsection{Development} % {{{

The Drupal project changed it's development workflow for the upcoming Drupal~8
release. Until then, the development was mostly driven through the issue
queues, which is a bug tracking system on \url{drupal.org}. While it is still
used for communication, bugs and new features, the development process features
a number of core initiatives which will stand for a major area of Drupal
\cite{DrupalInitiatives}. Each initiative has one or two initiative owners who
lead the development and coordinates the initiative while working closely with
Dries Buytaert.

Initiatives can be created by any Drupal developer, however Dries Buytaert
choses which initiative will be followed for the next major release of Drupal.
He also choses the initiative leaders and remains in close contact with them.
Furthermore each initiative has defined a small set of goals which should be
reached by the time of the next major release. For the Drupal~8 release, the
following initiatives will be followed.

\begin{itemize}
  \item Configuration Management
  \item Web Services Context Core Initiative
  \item Design 4 Drupal
  \item Drupal 8 Multi-Lingual Initiative
  \item HTML5
  \item Mobile
\end{itemize}

\begin{figure}[htbp]
  \centering
  \includegraphics[width=0.8\textwidth]{drupal/authors_by_month}
  \caption{Amount of distinct authors of Drupal Core over time. Due to their
  policy, only certain developers are allowed to commit to the repository.
  Furthermore the project only recently started adding author names to the
  commits.}
\end{figure}

% }}}

% }}}
